\section{Selecting an eight-call design from sparse calibration}
\label{app:sparse-pilot}
For the two attained designs in Section~\ref{sec:budget}, write $r_P(d)$
and $r_W(d)$ for single-question MSE, and let $F$ denote the centered
Williams-group subspace of dimension five. The precedence residual
space is $F\oplus G$, with $\dim G=6$. The projector identities in
Appendix~\ref{app:budgets} imply
\begin{equation}
 r_P(d)-r_W(d)=\frac{1}{22}
 \left(\frac{\|P_Gd\|^2}{6}-\frac{\|P_Fd\|^2}{5}\right).
 \label{eq:pilot-difference}
\end{equation}
Both terms admit unbiased eight-order estimates that remove every
additive position and pairwise precedence effect on each draw.

\paragraph{An explicit residual pilot.}
For each pair $A,B$ of the six Williams blocks, form the signed vector
$f=\mathbf1_A-\mathbf1_B$. These 15 vectors have squared norm eight.
Averaging $ff^\top/8$ uniformly gives $P_F/5$, by the sum of squared
pair differences among six group means. For $G$, use the six vectors
$g=\mathbf1_{+}-\mathbf1_{-}$ in Table~\ref{tab:pilot-directions}.
They are mutually orthogonal, have squared norm eight and annihilate
the constant, the position and precedence features, and all six
Williams-block indicators. They therefore span $G$. Uniformly
averaging $gg^\top/8$ gives $P_G/6$. These two identities hold for
every real response profile. They specialize to ternary differences
of binary correctness.

\begin{table}[h]
\centering
\caption{Sparse directions for the additional six-dimensional subspace.
Each tuple lists four display-order indices in the lexicographic convention of Appendix~\ref{app:budgets}.}
\label{tab:pilot-directions}
\begin{tabular}{ll}
\toprule
Positive coordinates & Negative coordinates \\
\midrule
$(0,7,16,23)$ & $(2,10,13,21)$ \\
$(0,7,17,22)$ & $(1,6,16,23)$ \\
$(1,6,17,22)$ & $(4,8,15,19)$ \\
$(2,11,13,20)$ & $(3,10,12,21)$ \\
$(3,11,12,20)$ & $(5,9,14,18)$ \\
$(4,9,14,19)$ & $(5,8,15,18)$ \\
\bottomrule
\end{tabular}
\end{table}

Sample either family with probability one half, then a uniform member.
Three eight-order supports appear in both families, giving 18 possible
blocks. Write $I_F(B),I_G(B)$ for the two membership indicators.
The probability of observing $B$ is
$q_B=I_F(B)/30+I_G(B)/12$. Given the eight observed $d_j$, use
\[
 Y_B=\frac{1}{176q_B}
 \left[\frac{I_G(B)(g_B^\top d)^2}{6}-
       \frac{I_F(B)(f_B^\top d)^2}{15}\right],
\]
with an absent direction contributing zero. Taking expectation over
$B$ recovers Eq.~\eqref{eq:pilot-difference}. Each order has inclusion
probability $1/3$, so the observed mean of the same eight responses is
also an unbiased estimate of the question's full-order gap. A separate
verifier checks the quadratic identity on every one- and two-coordinate
indicator profile and checks the signed vectors against the feature spans.

\paragraph{Using the precedence design for calibration.}
For the secondary rule draw $B$ from $P$ and let $p_{ab},w_{ab}$ be
the joint inclusion probabilities of distinct orders $a,b$ under
$P,W$. Every $p_{ab}$ is positive. The two covariance matrices have
zero row sums and equal diagonal entries. Expanding their difference
in squared pair differences gives the unbiased score
\[
 Z_B=\frac1{64}\sum_{\{a,b\}\subset B}
 \left(\frac{w_{ab}}{p_{ab}}-1\right)(d_a-d_b)^2 .
\]
This inverse-probability weighting uses the same 28 observed order
pairs \citep{horvitz1952}. The calibration-item means have exactly
the same risk as under fixed precedence sampling. On the three
previously published 3B panels the residual pilot had signed-score
standard deviations of 0.0155--0.0157 across items and pilot draws.
For the inclusion-weighted score the range was 0.0511--0.0613.
The primary pilot's calibration gap MSE exceeded fixed precedence
by 13.1--18.8\% on those development panels.

\paragraph{Risk with calibration charged.}
Let $C,H$ denote the fixed calibration and held-out rosters,
$Q$ the calibration sampling law and $\gamma$ the probability that
the calibration mean score is positive. Draw the selected deployment
sampler independently for each held-out question. Its conditional
mean gap error is zero for either choice. Hence its covariance with
the calibration gap error vanishes, even though the choice depends
on calibration observations. For $N=|C|+|H|$,
\begin{align}
 N^2\operatorname{MSE}(\widehat\Delta_{\rm policy})
 &= \sum_{i\in C}r_Q(d_i)
 +(1-\gamma)\sum_{i\in H}r_P(d_i)
 +\gamma\sum_{i\in H}r_W(d_i).
 \label{eq:charged-pilot}
\end{align}
The policy improves on always using $P$ precisely when the
expected held-out saving $\gamma\sum_{i\in H}(r_P-r_W)$ exceeds
the calibration cost $\sum_{i\in C}(r_Q-r_P)$. Both sides compare
the same total of eight calls per model and question.

\paragraph{Transfer protocol.}
The public registration fixed the four validation rosters, their
calibration/test splits and every sparse pilot block before GPU
inference. Within each benchmark we admitted distinct four-choice
questions, filtered explicit slot-relative answers and kept a single
question per normalized passage or standalone stem. All 24 numeric
prompts had to fit both tokenizers' 512-token limit. RACE had 135
eligible passages after tokenization, of which 128 were used. The
other three panels each contributed 256 questions. There were 294,
497 and 643 eligible metadata groups on ARC-Challenge, OpenBookQA and
Cosmos QA, respectively. No normalized question stem matched a stem
in the three 3B development panels. Earlier tests with a different
model pair shared 252 ARC-Challenge questions and 155 OpenBookQA
questions with this study. These four panels were outside development
of the eight-call rule. Exposure
to source data during model training is unknown.

A hash assigned one quarter of each roster to calibration. Each
method saw only its own eight paired correctness differences per
calibration item. We recorded its choice before computing any
held-out correctness profile. Full reference acquisition cost 43,008
model--item--order calls. Each separately evaluated policy would use
14,336 calls across the two models, including calibration. We compare
the registered choice with the exact held-out risks of both samplers.
We also select from the exact 24-order losses on those same calibration
items as a diagnostic costing three times as many pilot calls.

For each panel, 50,000 paired question resamples gave percentile
intervals for the held-out relative Williams gain. Simultaneous bounds
use the 0.625th and 99.375th percentiles per panel for a 95\% Bonferroni
family across all four. At fixed rosters, we rerandomized the pilot
20,000 times to estimate $\gamma$ in Eq.~\eqref{eq:charged-pilot}.
The remaining terms are exact. For a separate arithmetic check of the primary rule, we encode each calibration score distribution as an integer polynomial with coefficients counting sampler tickets. Multiplying the item polynomials gives the exact probability of a positive score sum. We use that probability for the primary policy MSE in Table~\ref{tab:pilot-charged}. Sampling variation across fresh
benchmarks is outside both calculations.

\begin{table}[h]
\centering
\small
\caption{Held-out benchmark diagnostics. MSEs use units of $10^{-6}$. The interval is the paired
Bonferroni interval for the relative W gain in percent. The last two columns show the prespecified
secondary choice and the choice obtained using all 24 orders on the calibration questions.}
\label{tab:pilot-holdout}
\begin{tabular}{@{}lrrrrcc@{}}
\toprule
Panel & P MSE & W MSE & Family interval, \% & $E_W/E_P$, \% & Secondary & 24-order pilot \\
\midrule
ARC-Challenge & 38.019 & 40.125 & $[-13.95,+3.41]$ & 47.97 & P & W \\
OpenBookQA & 34.294 & 37.130 & $[-18.25,+1.60]$ & 49.21 & W & P \\
Cosmos QA & 33.720 & 34.737 & $[-12.85,+6.42]$ & 46.83 & P & P \\
RACE & 38.430 & 41.594 & $[-22.33,+6.06]$ & 49.20 & W & P \\
\bottomrule
\end{tabular}
\end{table}

\begin{table}[h]
\centering
\small
\caption{MSE of the full-panel gap in units of $10^{-6}$, with eight calls per model per question for every policy.
Primary policy MSE integrates over all possible pilot order assignments. Secondary policy MSE uses the registered
20,000 randomizations. The last column gives the primary pilot's percent change in MSE on the calibration questions
alone, relative to using P on those same questions.}
\label{tab:pilot-charged}
\begin{tabular}{@{}lrrrrr@{}}
\toprule
Panel & Always P & Always W & Residual pilot & Precedence pilot & Calibration MSE change, \% \\
\midrule
ARC-Challenge & 28.100 & 29.204 & 29.856 & 28.700 & $+17.26$ \\
OpenBookQA & 26.178 & 28.091 & 28.112 & 26.942 & $+17.88$ \\
Cosmos QA & 26.539 & 27.233 & 27.734 & 26.823 & $+12.20$ \\
RACE & 29.766 & 32.679 & 31.512 & 30.591 & $+13.51$ \\
\bottomrule
\end{tabular}
\end{table}

The secondary rule chose W on OpenBookQA and RACE and P on ARC-Challenge
and Cosmos QA, matching the held-out lower-risk design twice. Its
expected full-panel MSE was $1.07$--$2.92\%$ higher than fixed P
across the four benchmarks. The exact 24-order calibration diagnostic
selected W on ARC-Challenge, where P had lower held-out risk. This
case illustrates a source of uncertainty present even when calibration
order effects are fully observed. Here, changing the questions changed the ranking. Table~\ref{tab:pilot-charged}
separates the higher risk of the primary sampler on calibration
questions from errors in selecting the held-out sampler.
An independent reconstruction using the publisher keys and the raw
numeric choices matched every released correctness-difference mask.

